Det är upp till varje RU att ansvara för sin egen IT-lösning för C-DAS\@.
Att lokförare följer C-DAS-hastighetsrekommendationer är kritiskt för att den nya paradigmen med TMS/C-DAS ska fungera.
%Forskning på tillämpningen av C-DAS ur lokförares perspektiv är ringa.
\begin{enumerate}
    \item Hur ser benägenheten ut för lokförare att följa hastighetsrekommendationer från en DAS-applikation idag?
    \item Vad bör utvecklare av ett C-DAS-ombordsystem beakta ur ett lokförarperspektiv (för vagt kanske)?
    \begin{itemize}
        \item Vad är ett bra sätt att återge hastighetsrekommendationer från C-DAS-ombordsystem för lokförare?
        \item Vad är ett lämpligt sätt att tillämpa hastighetsrekommendationer för att öka lokförares benägenhet att följa dem?
    \end{itemize}
    \item Av vilka data består kommunikation från en C-DAS-lösning som har hög benägenhet att efterlevas av lokförare?
    \item På vilket sätt skiljer sig behoven åt mellan örare av godståg respektive persontåg avseende en C-DAS-lösning?
\end{enumerate}
